% !TEX TS-program = lualatex

\documentclass[border = 3pt]{standalone}

%%%
% Des couleurs nommées faciles d'emploi via le package \xcolor.
%%%
\usepackage[svgnames]{xcolor}

%%%
% La bibliothèque \luadraw allie une facilité d’utilisation
% à un rendu particulièrement soigné.
%%%
\usepackage[3d]{luadraw}

\begin{document}

\begin{luadraw*}{name = for-stacked-cubes}
local ld = luadraw

local Origin = ld.pt3d.Origin
local M      = ld.pt3d.M

local vecI = ld.pt3d.vecI
local vecJ = ld.pt3d.vecJ
local vecK = ld.pt3d.vecK

------
-- Préréglages.
------
local h_max = 7
local pal  = {"Teal", "LightBlue"}

------
-- ¨Def de la zone graphique.
------
local xy_sup = h_max
local z_sup  = h_max*(h_max + 1) / 2

local graphview = ld.graph3d:new{
  window3d = {0, xy_sup, 0, xy_sup, 0, z_sup},
  viewdir  = {155, 70},
  adjust2d = true,
  bbox     = false
}

------
-- Réglages liés aux lignes.
------
graphview:Linewidth(6)

------
-- Empilons des cubes.
------
local z_pos   = 0
local i_col   = 1
local nb_cols = #pal

for a = h_max, 1, -1 do
  graphview:Dpoly(
    ld.parallelep(
      M(0, 0, z_pos),
      a*vecI, a*vecJ, a*vecK
    ),
    {
      mode  = ld.mFlat,
      color = pal[i_col]
    }
  )

  z_pos = z_pos + a
  i_col = i_col % nb_cols + 1
end

------
-- Montrons le résultat de notre oeuvre.
------
graphview:Show()
\end{luadraw*}

\end{document}
